\clearpage\thispagestyle{empty}
\pdfbookmark[0]{Quellen zur Gerling-Besprechung}{gerling1812-sources}
\begin{center}\textbf{Quellen zur Gerling-Besprechung}\end{center}
\noindent Diese KI-Übersetzung ins heutige Deutsch folgt Carl Friedrich Gauß, \textit{Werke}, Bd.~VI (Göttingen, 1874), S.~542--543. Die Besprechung erschien zuerst unter Göttingen in den \textit{Göttingischen gelehrten Anzeigen}, Stück~190, 28. November 1812, S.~1889--1891.
\par\smallskip\noindent Im Erstdruck verläuft die Senkrechte vom Mittelpunkt der Erde zum Mittelpunkt des bedeckten oder verfinsterten Gestirns. Die Werkausgabe nennt nur den Mittelpunkt des bedeckten oder verfinsterten Gestirns. Die Besprechung nennt 46 Seiten und I.~C.~Baier. Der \href{https://ci.nii.ac.jp/ncid/BD21343055}{Katalogeintrag zum Exemplar der Universität Hokkaido} nennt 48 Seiten, eine ungezählte gefaltete Tafel und J.~C.~Baier.
\par\smallskip\noindent \href{https://gdz.sub.uni-goettingen.de/id/PPN236007467}{Text der Werkausgabe}. Erstdruck, \href{https://images.p.sub.uni-goettingen.de/iiif/image/gdz:PPN319721507_1812:00001933/full/1650,/0/default.jpg}{S.~1889}, \href{https://images.p.sub.uni-goettingen.de/iiif/image/gdz:PPN319721507_1812:00001934/full/1650,/0/default.jpg}{S.~1890} und \href{https://images.p.sub.uni-goettingen.de/iiif/image/gdz:PPN319721507_1812:00001935/full/1650,/0/default.jpg}{S.~1891}.
\par\smallskip\noindent Übersetzung und Satz wurden von KI erstellt.
\clearpage
